\chapter{State of the Art}
\label{chapter:state_of_the_art}

\section{Current Approaches in Flood Modeling}

\subsection{Overview of Flood Simulation Paradigms}

Flood modeling has changed a lot in the ten years. It used to be about looking at data to make predictions. Now it is about making 
predictions in time. This evolution is driven by the increasing frequency of extreme meteorological events, such as the catastrophic 
DANA experienced in the Spanish Mediterranean in October 2024. People wrote about this in papers by \cite{SeguraBeltran2025} and 
\cite{Muoz2025}. There are three ways that people make flood models now. The first way is to use based hydrodynamic models. The second 
way is to use computational models, such, as Cellular Automata. The third way is to use something called Digital Twins, which's a new 
idea that is part of something bigger called Society 5.0. Digital Twins are a new way to do it.


\subsection{Hydrodynamic Modeling and Shallow Water Equations (SWE)}

The traditional gold standard in flood simulation uses the Shallow Water Equations (SWE) to model floods. These equations come from 
the Navier-Stokes equations. Assume that the horizontal length is much bigger than the vertical depth.

Advanced methods like Godunov-type schemes are often used to solve these equations on meshes. As noted by \cite{Aleksyuk2017}, these 
schemes are important for handling terrain with areas where water gets shallow and the bottom changes suddenly. In areas the water 
flow can change from slow to fast and vice versa. The models use solvers to capture shock waves, which is crucial for simulating dam 
breaks or violent flash floods in steep areas. 

Recently these models have been used in places like the Northern Caucasus as discussed by \cite{Vasileva2021}. They work with automated 
weather and water monitoring systems. These systems allow for simulations in small river areas where flash floods happen quickly. 
However these detailed models are very costly to run and require computers or special graphics cards to work fast enough for large 
areas or in real-time. The high cost is a problem, for using them widely as noted by \cite{Aleksyuk2017}.


\subsection{Discrete Models: The Rise of Cellular Automata (CA)}

As an alternative to the heavy computational requirements of differential equation solvers, Cellular Automata (CA) have 
emerged as a robust framework for flood inundation modeling. Unlike continuous models, CA operate on a discrete grid 
where each cell changes state based on local transition rules.

A significant milestone in this field was the development of the multilayer cellular automata ($m:n-CA^k$) by \cite{Fonseca2011}. 
Originally applied to slab avalanche configurations, this model proved that complex geophysical phenomena 
could be effectively represented by separating the geographical data into distinct layers (GIS-based) and applying formal 
transition rules (represented via SDL - Specification and Description Language). This "Data-as-Interface" strategy allows 
for a clear decoupling between the model logic and the implementation, facilitating extensibility.

Current research has further refined CA-based solvers to include dynamic-wave approximations, achieving 
accuracies comparable to traditional Finite Volume methods but with execution speeds up to 300% faster. Models like CA-ffé 
(Cellular Automata fast flood evaluation) have demonstrated the ability to predict maximum inundation depths in seconds, making 
them ideal for exploratory water system planning and real-time emergency response \parencite{Jamali2019}.


\subsection{Digital Twins and the Society 5.0 Framework}

The most recent frontier in flood management is the integration of simulation models into Digital Twin (DT) architectures. 
Within the vision of Society 5.0, a Digital Twin is not merely a 3D visualization but a dynamic, real-time replica of the 
physical environment that ingests data from IoT sensors, satellite imagery, and meteorological radars \parencite{FonsecaiCasas2026}.

According to \cite{FonsecaiCasas2026}, Society 5.0 leverages Industry 4.0 technologies—such as Big Data, AI, and Digital 
Twins—not just for industrial efficiency but for human-centric social benefit, specifically in disaster resilience. In this 
context, a flood model becomes an "Open Model" that supports decision-making by allowing authorities to simulate "what-if" 
scenarios in real-time. For instance, during the 2024 DANA event, the lack of integrated real-time modeling was cited as a 
major hurdle in effective emergency management. The transition toward "Digital Twin Smart Cities" (DTSC) 
aims to bridge this gap by providing a continuous synchronization between the physical river basin and its digital counterpart 
\parencite{Mazzetto2024}.


\subsubsection{Distributed Communication in Early Warning Systems (MQTT)}

The transition from static simulations to reactive Digital Twins requires a communication backbone capable of handling high-frequency 
data exchange with minimal overhead. In the context of flood monitoring, the Message Queuing Telemetry Transport (MQTT) protocol has 
emerged as a standard for real-time environmental telemetry \parencite{Dionsio2025}. Unlike traditional request-response models (HTTP), 
MQTT's publish-subscribe architecture allows for a decoupled interaction between the simulation engine and the visualization layer. 
This is particularly critical in flash flood scenarios where sub-second latency in data dissemination can significantly impact the 
efficacy of emergency response.


\subsubsection{Software Design Patterns for Scientific Modeling (Hexagonal Architecture)}

A common pitfall in scientific software development is the "Big Ball of Mud", where physical logic is tightly coupled with data I/O and 
specific library dependencies (e.g., GIS formats or networking protocols). To ensure the longevity and testability of the simulator, 
this research adopts the Hexagonal Architecture (also known as Ports and Adapters), \cite{Cockburn2005}; \cite{Nunkesser2022}. By isolating the Cellular Automata Core 
(the "Inside") from the infrastructure concerns (the "Outside"), such as HDF5 file readers or MQTT publishers, the system achieves a 
high degree of technological agnosticism. This architectural decoupling allows for the substitution of data sources or visualization 
engines without modifying the underlying hydraulic transition rules. In modern scientific computing, this approach is increasingly 
recognized as a best practice to manage the complexity of multi-layered models ($m:n-CA^k$) and to facilitate the integration of 
heterogeneous geospatial data.


\subsection{Data-Driven and Empirical Approaches}

Complementing the numerical models are empirical approaches focused on peak flow estimation and hydrological response. 
\cite{Muoz2025} highlight the importance of "reconstructing" flood events using post-event field data and empirical 
formulas (such as the Manning equation and the rational method) to validate hydraulic models. Furthermore, the use of 
public datasets (AEMET, SAIH, Copernicus) has become standard for pre-processing inputs, although the "spatialization" 
of rainfall remains a challenge in small, ungauged catchments \parencite{Vasileva2021}.


\section{Analysis of Shortfalls in Existing Solutions}

The big improvements in modeling and discrete simulation that we talked about earlier do not change the fact that the really bad 
flash floods in 2024 like the DANA in the Júcar River Basin showed us that there are big problems with the way things are done now. 
These problems can be put into four groups: it takes a long time to get the results, the system is not flexible, the data is not shared 
well, and the flash flood systems do not have a simple way for people to make decisions. The flash flood systems have these problems 
because of the latency, the architectural rigidity, the data integration silos and the lack of a human-centric decision-making 
interface, in the flash flood systems.


\subsection{Computational Latency vs. Flash Flood Dynamics}

The primary shortfall of physically-based models (SWE solvers) remains their execution time. Achieving high-precision results on 
complex, unstructured meshes requires high-performance computing resources that are often unavailable or too slow to deploy during 
the narrow "lead time" of a flash flood. In events where the time from peak rainfall to peak flow is measured in minutes rather than 
hours, the computational overhead of traditional numerical methods prevents them from being used as proactive forecasting tools. 
Instead, they are frequently relegated to post-event forensic analysis \parencite{Muoz2025}.


\subsection{Monolithic Architectural Rigidity}

Many flood simulation platforms are built as one system. This makes it very hard to add features or change parts of it. For example, 
you can't just change the way water soaks into the ground or add a weather data source without overhauling everything. This lack of 
flexibility makes it tough to create models like those, in the Society 5.0 plan \parencite{FonsecaiCasas2026}. These systems also only work with specific data formats.
That makes it hard for them to work with GIS tools and real-time simulation engines.


\subsection{Simplified Representation of Spatial Heterogeneity}

While standard Cellular Automata (CA) offer the speed required for real-time applications, they often suffer from an 
oversimplification of the physical environment. Most CA implementations treat the terrain as a single-dimensional grid, 
failing to represent the complex interactions between different layers of information—such as soil saturation, 
subterranean drainage systems, and varying urban roughness—simultaneously. The inability to handle "multilayered" spatial 
data in a unified tensor-like structure limits the model's capacity to simulate the non-linear behavior of water in 
urbanized basins \parencite{Fonseca2011}.


\subsection{The Socio-Technical Gap in Emergency Management}

A final, critical shortfall is the disconnect between the technical output of simulation models and the operational needs 
of emergency responders. As highlighted by regarding the 2024 DANA crisis, the lack of a dynamic, real-time 
"Digital Twin" meant that decision-makers were operating with fragmented and delayed information. Existing models do not 
easily integrate with the modern IoT and radar infrastructure needed to provide a continuous, synchronized digital 
representation of the evolving disaster.


\section{Problem Statement}

The main problem that this thesis is trying to fix is that current flood modeling systems are not good enough. They are not fast, 
accurate, and flexible all at the same time. This means they are not very useful for making decisions during really bad flash floods.

There is a gap between two types of models. The first type, hydrodynamic solvers, is very precise but very slow. The second type, 
discrete models, is fast but not very precise. This gap is made worse because we do not have a way to handle all the complex data 
that we have about the earth. We need a system that can handle all this data in a way that's scalable.

The 2024 DANA event showed us that our old way of predicting floods is not good enough for what society needs today. We need a way 
of doing things. We need systems that can take in data in time, process it quickly ,and give us results that we can use to make 
decisions. These systems are called Reactive Digital Twins.

So this research is trying to solve some technical problems. These problems are:

\begin{enumerate}
    \item \textbf{Modeling Complexity}: How to create a model that can show us how different things like rain and water flowing over 
    the ground interact with each other. We want to use a cellular automata to do this.
    \item \textbf{Architectural Decoupling}: How to design a software architecture (Hexagonal/Ports and Adapters) that 
    allows the simulation engine to remain agnostic to specific GIS data formats and storage backends.
    \item \textbf{Performance and Scalability}: How to implement this logic in a high-performance environment (C++) that 
    leverages modern data structures (Tensors) to ensure the simulation speed remains well below the real-time threshold.
\end{enumerate}

If we can solve these problems then we can create a way of managing floods. This new way will be able to take in lots of data, about 
the environment and turn it into decisions that can save lives. Flood management tools will be much better. They will be able to help 
us make decisions during really bad floods.
